\section{\label{sec:4} Discussion}
Earmolds repeatedly disrupted vertical localization, but participants relearned sound localization with both mold sets within successive five-day adaptation periods. Individual adaptation performance did not improve with the second earmolds and learning a second set did not interfere with the previously learned representation. Localization accuracy with free ears slightly decreased throughout the experiment, whereas elevation perception remained unaffected. Spectral information of free ears correlated with behavior, and was similarly reduced by both mold sets. Acoustic dissimilarities between successive pinna shapes (free ears vs. the first molds and the first vs. the second molds) were similar, though each set induced changes at distinct frequencies.

\subsection{Localization with free ears} 